\section*{Overview}

Bitset optimization is the habit of rephrasing boolean-state transitions so that one machine word updates 64 states at
once. In competitive programming this often turns borderline \(O(n^2)\) boolean DP or dense-graph processing into
something that actually passes.

\section*{When to Use It}

It is especially effective when:

\begin{itemize}
  \item the state is yes/no rather than numeric,
  \item the transition is based on shifts, AND, OR, XOR, or intersection,
  \item the dimension is large enough that a \(64\times\) constant-factor improvement matters.
\end{itemize}

\section*{Core Idea}

Store a boolean array as packed bits instead of \texttt{bool}s. Then operations such as
\[
\texttt{reachable |= reachable << w}
\]
update many states at once.

\section*{Key Insight}

The optimization is not magic. It works because the CPU already knows how to shift, mask, and combine entire words very
quickly. If the transition can be written in those operations, you get parallelism ``for free'' without threads or SIMD
intrinsics.

\section*{Worked Problem}

\paragraph{Problem.}
Given positive integers \(a_1, a_2, \dots, a_n\) with total sum at most \(S\), determine which subset sums are
reachable.

\paragraph{Why bitsets fit.}
The ordinary DP is
\[
\texttt{dp[s] = dp[s] or dp[s - a_i]}.
\]
With a bitset, one item update becomes \texttt{dp |= dp << a\_i}. That replaces \(O(S)\) scalar transitions with
\(O(S / 64)\) word operations.

\section*{Problem Pattern}

The same trick appears in:

\begin{itemize}
  \item subset-sum and knapsack reachability,
  \item dense graph transitive closure with row bitsets,
  \item maintaining large character sets or forbidden masks,
  \item fast set intersections in Mo-style or offline problems.
\end{itemize}

\section*{Correctness Intuition}

The packed representation does not change the DP semantics. Bit \(s\) is still ``sum \(s\) is reachable.'' The bit
operations merely update 64 such booleans in parallel.

\section*{Complexity Analysis}

If the original boolean DP is \(O(NM)\), the bitset version is usually \(O(NM / W)\), where \(W\) is the machine word
size, typically \(64\).

\section*{Implementation}

The sample code gives a reusable subset-sum bitset wrapper. In actual contest code I often use \texttt{std::bitset} for
compile-time bounds and a manual \texttt{vector<uint64\_t>} wrapper when the bound is only known at runtime.

\section*{Common Pitfalls}

\begin{itemize}
  \item Applying bitsets to numeric DP where each state stores a large value, not a boolean.
  \item Forgetting the memory cost; a bitset of size \(10^7\) is already noticeable.
  \item Using \texttt{std::bitset} when the size is only known at runtime.
  \item Expecting asymptotic improvements when the transition itself is not bit-parallel.
\end{itemize}

\section*{Variants / Extensions}

\begin{itemize}
  \item Manual dynamic bitsets with \texttt{uint64\_t}.
  \item Bitset adjacency rows for dense graph BFS or transitive closure.
  \item Using \texttt{\_\_builtin\_popcountll} on words to count set bits quickly after intersections.
\end{itemize}

\section*{Practice Problems}

\begin{itemize}
  \item Subset-sum reachability with large total sum.
  \item Dense graph reachability or clique-style adjacency intersections.
  \item Any boolean DP whose transition is shift-or-mask based.
\end{itemize}

\section*{References}

\begin{itemize}
  \item \href{https://usaco.guide/plat/bitsets}{USACO Guide: Bitsets}.
  \item \href{https://codeforces.com/catalog}{Codeforces Catalog}.
  \item \href{https://github.com/kth-competitive-programming/kactl}{KACTL}.
\end{itemize}
